%%=============================================================================
%% Proof-of-Concept: React-applicatie met Application Insights
%%=============================================================================

\chapter{\IfLanguageName{dutch}{Proof-of-Concept: React-applicatie met Application Insights}{Proof-of-Concept: React Application with Application Insights}}%
\label{ch:poc}

\section{\IfLanguageName{dutch}{Opzet van de applicatie}{Application Setup}}

Als basis voor dit onderzoek werd een proof-of-concept React-applicatie ontwikkeld die drie meerstapsformulieren bevat. De applicatie werd gebouwd met als doel het effect van usability-testen op verschillende momenten in het ontwikkelingsproces te kunnen meten aan de hand van echte gebruikersinteracties. De broncode is beschikbaar via GitHub\footnote{\url{https://github.com/fienteugels/fien-teugels-BAP}}.

De applicatie bestaat uit drie onafhankelijke formulieren, elk met een eigen thema en doelgroep:

\begin{itemize}
  \item \textbf{Form A — Reisaanvraag:} een meerstapsformulier waarmee gebruikers een reisaanvraag kunnen indienen. Het formulier omvat negen stappen en bevat conditionele logica op basis van het gekozen reistype, continent en vervoersmiddel.
  \item \textbf{Form B — Evenementregistratie:} een meerstapsformulier voor het registreren van een evenement. Het formulier telt tien stappen en past de vragen aan op basis van het gekozen evenementtype (festival, conferentie, sportevenement, \ldots).
  \item \textbf{Form C — Verblijfsaanvraag:} een meerstapsformulier voor een verblijfsaanvraag. Het formulier telt tien stappen en 65 invulvelden, waarvan een deel conditioneel is op basis van het gekozen verblijfstype.
\end{itemize}

Elk formulier is bereikbaar via een afzonderlijke route binnen de applicatie en maakt gebruik van een gedeelde formuliercomponent.

\section{\IfLanguageName{dutch}{Formulierarchitectuur}{Form Architecture}}

De drie formulieren zijn gebouwd bovenop een herbruikbare \texttt{FormShell}-component. Deze component beheert de staptoestand, de navigatielogica (volgende stap, vorige stap, validatie) en de verzending van telemetry-events naar Application Insights. Door de logica te centraliseren in één component wordt consistentie gegarandeerd over alle formulieren heen en kan telemetry uniform worden bijgehouden.

Elk formulier geeft een uniek \texttt{formId} mee aan de \texttt{FormShell}-component. Dit prefix wordt gebruikt bij alle gegenereerde Application Insights-events, zodat de data per formulier kan worden gefilterd. De stappen van een formulier worden gedefinieerd als een array van componenten, waarbij elke stap zijn eigen validatieregels en conditionele weergavelogica bevat.

Conditionele velden worden getoond of verborgen op basis van eerder ingegeven antwoorden via de \texttt{useState}- en \texttt{useEffect}-hooks van React. Wanneer een antwoord een tegenstrijdige combinatie vormt met een eerder antwoord, verschijnt een inline waarschuwing zonder de gebruiker te blokkeren.

\section{\IfLanguageName{dutch}{Integratie met Microsoft Application Insights}{Integration with Microsoft Application Insights}}

Elke gebruikersinteractie die relevant is voor de usability-analyse wordt automatisch doorgestuurd naar Microsoft Application Insights via de \texttt{@microsoft/applicationinsights-web} SDK. Application Insights fungeert als de centrale dataverzamelplaats voor alle telemetry en maakt het mogelijk om via Kusto Query Language (KQL) gerichte analyses uit te voeren op de verzamelde data.

\subsection{\IfLanguageName{dutch}{Geregistreerde events}{Tracked Events}}

Tabel~\ref{tab:ai-events} geeft een overzicht van alle custom events die door de applicatie worden bijgehouden, met hun naamgeving op basis van het \texttt{formId}-prefix.

\begin{table}[h]
  \centering
  \caption[Custom Application Insights events]{Overzicht van de custom events per formulier}
  \label{tab:ai-events}
  \begin{tabular}{lp{8cm}}
    \hline
    \textbf{Event (prefix\_Naam)} & \textbf{Beschrijving} \\
    \hline
    \texttt{FormX\_StepView}             & Gebruiker bekijkt een stap; registreert het stapnummer en de sessie-ID. \\
    \texttt{FormX\_StepComplete}         & Gebruiker voltooit een stap succesvol; registreert de doorlooptijd in milliseconden. \\
    \texttt{FormX\_StepBack}             & Gebruiker navigeert terug naar een vorige stap. \\
    \texttt{FormX\_StepValidationFailed} & Validatie mislukt op een stap; registreert het stapnummer en het type fout. \\
    \texttt{FormX\_FieldChange}          & Gebruiker wijzigt een veldwaarde; registreert de veldnaam. \\
    \texttt{FormX\_DropOff}             & Gebruiker verlaat het formulier zonder het te voltooien; registreert de laatste stap. \\
    \texttt{FormX\_FormSubmitted}        & Gebruiker dient het formulier succesvol in. \\
    \texttt{FormX\_FeedbackSubmitted}    & Gebruiker dient het in-app feedbackformulier in; bevat rating en NPS-score. \\
    \hline
  \end{tabular}
\end{table}

Waarbij \texttt{FormX} het prefix is van het betreffende formulier, bijvoorbeeld \texttt{FormA}, \texttt{FormB} of \texttt{FormC}.

De events bevatten naast de naam ook een \texttt{session\_Id} waarmee alle interacties van één gebruiker kunnen worden gekoppeld. Dit maakt het mogelijk om volledige gebruikerssessies te reconstrueren en funnel-analyses uit te voeren per sessie in plaats van per losse paginaweergave.

\subsection{\IfLanguageName{dutch}{Kusto-queries voor analyse}{Kusto Queries for Analysis}}

De telemetry-data werd geanalyseerd via KQL-queries in de Azure Portal (Logboeken). De volgende queries werden gebruikt om de drie kernmetrics te berekenen:

\begin{itemize}
  \item \textbf{Voltooiingsratio en funnel:} telt het aantal unieke sessies per stap op basis van \texttt{StepComplete}-events, waarna de drop-off per stap zichtbaar wordt.
  \item \textbf{Gemiddelde stap-doorlooptijd:} berekent per stap de gemiddelde duur in milliseconden op basis van de \texttt{duration\_ms}-eigenschap van \texttt{StepComplete}-events.
  \item \textbf{Validatiefouten per sessie:} telt het aantal \texttt{StepValidationFailed}-events per sessie en berekent het gemiddelde over alle sessies.
\end{itemize}

\section{\IfLanguageName{dutch}{In-app feedbackformulier}{In-app Feedback Form}}

Na het indienen van elk formulier krijgt de gebruiker automatisch een kort feedbackformulier te zien. Dit formulier bevraagt:

\begin{itemize}
  \item een algemene beoordeling op een schaal van 1 tot 5;
  \item de ervaren moeite bij het invullen;
  \item de geschatte invultijd;
  \item de duidelijkheid van de stapindicatie;
  \item de beoordeling van waarschuwingen bij tegenstrijdige combinaties;
  \item de beoordeling van de conditionele velden;
  \item de vlottheid van de navigatie;
  \item open vragen over verwarring en verbeterpunten;
  \item een NPS-score (0–10).
\end{itemize}

De antwoorden worden via het \texttt{FormX\_FeedbackSubmitted}-event als JSON opgeslagen in Application Insights en zijn opvraagbaar via dezelfde KQL-interface. Dit zorgt ervoor dat kwantitatieve telemetry en subjectieve feedback aan dezelfde sessie kunnen worden gekoppeld.
